\chapter{Tecnologías de Reciclaje y Gestión de Residuos}
\label{cap:reciclaje}

El cierre del ciclo de vida de los materiales plásticos es un requisito técnico indispensable en la ingeniería de procesos contemporánea. Las tecnologías de reciclaje se clasifican en cuatro niveles de gestión y valorización de residuos:

\section{Reciclaje Mecánico}

Corresponde a la valorización física y es la técnica industrial más extendida comercialmente. Abarca procesos de clasificación ( NIR), molienda, lavado intensivo y re-extrusión del desecho posconsumo.

\subsection{Degradación en el Reciclaje Mecánico}
Durante los ciclos sucesivos de procesamiento por extrusión de plásticos como el PET o el PP, el material se expone a altos esfuerzos cortantes de cizalla mecánica y elevadas temperaturas en presencia de trazas de oxígeno molecular. Esto desencadena reacciones de termooxidación y escisión de cadena (\textit{chain scission}), provocando una disminución sistemática del peso molecular promedio y una consecuente reducción drástica de la viscosidad en fundido y de las propiedades de resistencia mecánica a la tracción y resistencia al impacto de la resina reciclada.

\subsection{Incompatibilidad Termodinámica de Mezclas}
El reciclaje mecánico de residuos plásticos mezclados es un desafío debido a que casi todos los polímeros diferentes son termodinámicamente inmiscibles entre sí debido a su baja entropía de mezcla ($\Delta G_m > 0$). Una mezcla fundida de PE y PP forma fases separadas físicamente con una interfase débil que carece de adhesión mecánica molecular. Para solucionar este problema se emplean compatibilizantes (copolímeros en bloque donde cada bloque es afín a uno de los plásticos de la mezcla), reduciendo la tensión interfacial ($\gamma$):
\begin{equation}
    \gamma = \gamma_0 - R T \Gamma
    \label{eq:compatibilizacion}
\end{equation}

\section{Reciclaje Químico (Valorización Terciaria)}

El reciclaje químico transforma las cadenas poliméricas de desecho en compuestos químicos de bajo peso molecular o en sus respectivos monómeros originales purificados:

\begin{description}
    \item[Despolimerización por Solvólisis:] Aplicable a polímeros obtenidos por pasos (como el PET, nailon, poliuretanos). Utiliza solventes reactivos calientes como etilenglicol (glicólisis) o metanol (metanólisis) en presencia de catalizadores de transesterificación para deshacer reversiblemente los enlaces éster o amida.
    \item[Pirólisis y Craqueo Térmico:] Descomposición térmica controlada a altas temperaturas en atmósfera libre de oxígeno de poliolefinas de desecho para obtener una mezcla de hidrocarburos llamada nafta pirolítica, la cual se reintroduce en los procesos de refinamiento de petróleo para producir resinas plásticas vírgenes.
\end{description}
